\section*{Capítulo \ref{ch:g7:geological-time} --- Las edades de la Tierra}

\begin{solution}{exo:g7:geological-time:1}
(a) \omterm{def:g7:geological-time:unconformity}{discordancia}; (b) inclusión; (c) superposición; (d) sucesión
faunística; (e) intersección.
\end{solution}

\begin{solution}{exo:g7:geological-time:2}
(a) La arenisca roja, abajo, es la más antigua; la arenisca amarilla,
arriba, es la más reciente. (b) El \omterm{def:g7:geological-time:superposition}{principio de superposición}. (c) Las
capas no se han dado la vuelta desde que se formaron.
\end{solution}

\begin{solution}{exo:g7:geological-time:3}
(a) Paleozoico, Mesozoico, Cenozoico. (b) El Cenozoico. (c) El Mesozoico.
(d) El Jurásico viene justo después del Triásico; el Cretácico viene justo
antes del Paleógeno.
\end{solution}

\begin{solution}{exo:g7:geological-time:4}
(a) Falso: la \omterm{def:g7:geological-time:relative-dating}{datación relativa} solo ordena las \omterm{def:g3:stones-pebbles-sand:rock}{rocas}; una edad en años es
una \omterm{def:g7:geological-time:relative-dating}{edad absoluta}. (b) Verdadero. (c) Falso: los \omterm{def:g3:stones-pebbles-sand:pebble}{cantos} son más antiguos
que el conglomerado que los contiene (\omterm{def:g7:geological-time:inclusion}{principio de inclusión}).
(d) Verdadero (\omterm{def:g7:geological-time:fossil-succession}{principio de sucesión faunística}).
\end{solution}

\begin{solution}{exo:g7:geological-time:5}
(a) Más reciente: corta la capa 3, y una \omterm{def:g3:stones-pebbles-sand:rock}{roca} que corta a otra es más
reciente (intersección). (b) Más antiguo: la capa 4 descansa encima sin
estar cortada. (c) Capa 1, capa 2, capa 3, el granito, capa 4. (d) La
\omterm{def:g5:water-shapes-the-land:erosion}{erosión} desgastó las \omterm{def:g3:stones-pebbles-sand:rock}{rocas} que tenía encima hasta que el granito llegó a
la superficie; después, la capa 4 se depositó sobre esa superficie de
\omterm{def:g5:water-shapes-the-land:erosion}{erosión}, una \omterm{def:g7:geological-time:unconformity}{discordancia}.
\end{solution}

\begin{solution}{exo:g7:geological-time:6}
(a) $538.8 - 251.9 = 286.9$ millones de años. (b) $251.9 - 66 = 185.9$
millones de años. (c) 66 millones de años. (d) El Paleozoico es la más
larga, y el Cenozoico, la más corta (por ahora).
\end{solution}

\begin{solution}{exo:g7:geological-time:7}
(a) $4540 - 538.8 = 4001.2$ millones de años. (b) $4001.2 \div 4540
\approx 0.881$: alrededor del \textbf{88~\%}. (c) Antes del Paleozoico, los
animales con concha todavía no eran frecuentes, y los seres vivos blandos
casi nunca se convierten en \omterm{def:g7:sediments-and-fossils:fossil}{fósiles}; las conchas y los huesos, que
fosilizan bien, solo son frecuentes a partir del Paleozoico.
\end{solution}

\begin{solution}{exo:g7:geological-time:8}
(a) $4540 \div 100 = 45.4$ millones de años por metro. (b) $538.8 \div 45.4
\approx 11.9$: a unos \qty{11.9}{m} de la meta. (c) $66 \div
45.4 \approx 1.45$: a unos \qty{1.45}{m}. (d) $0.3 \div 45.4 \approx
0.0066$~m, es decir, unos \qty{0.66}{cm}, menos de un centímetro: un
corredor apenas lo notaría.
\end{solution}

\begin{solution}{exo:g7:geological-time:9}
(a) El Silúrico fue primero. Sí: la \omterm{def:g3:stones-pebbles-sand:rock}{roca} más antigua está debajo de la más
reciente. (b) Las capas grises se depositan horizontales; se ponen de
canto; se desgastan; la \omterm{def:g3:stones-pebbles-sand:sand}{arena} roja se deposita encima. (c) Es una
superficie de \omterm{def:g5:water-shapes-the-land:erosion}{erosión} entre \omterm{def:g3:stones-pebbles-sand:rock}{rocas} de edades muy distintas: el largo tiempo
durante el cual las capas grises se levantaron y se desgastaron no dejó
ninguna capa.
\end{solution}

\begin{solution}{exo:g7:geological-time:10}
(a) Las dos lutitas con amonites tienen la misma edad, y también las dos
calizas con erizos de \omterm{def:g2:water-on-earth:ocean}{mar}: mismos \omterm{def:g7:sediments-and-fossils:fossil}{fósiles}, misma edad (\omterm{def:g7:geological-time:fossil-succession}{principio de
sucesión faunística}). (b) La más antigua es la arenisca con trilobites del
acantilado A (está debajo de la lutita con amonites, que tiene la misma
edad que la capa de abajo del acantilado B). La más reciente es la arcilla
con hojas del acantilado B. (c) De abajo arriba: arenisca con trilobites;
lutita con amonites; caliza con erizos de \omterm{def:g2:water-on-earth:ocean}{mar}; arcilla con hojas.
\end{solution}

\begin{solution}{exo:g7:geological-time:11}
(a) El esquisto. Está debajo del conglomerado (superposición), y los
\omterm{def:g3:stones-pebbles-sand:pebble}{cantos} del conglomerado son trozos de esquisto (inclusión); además, la
superficie de \omterm{def:g5:water-shapes-the-land:erosion}{erosión} corta el esquisto. (b) Después: la arenisca se
depositó horizontal, así que se inclinó más tarde. (c) Anterior: la arcilla
cubre la \omterm{def:g7:earthquakes:fault}{falla} sin estar rota. (d) Se forma el esquisto; se desgasta
(superficie ondulada); se depositan el conglomerado, la arenisca y la
caliza, en ese orden; las capas se inclinan y la \omterm{def:g7:earthquakes:fault}{falla} las rompe --- el
corte por sí solo no dice cuál de las dos cosas fue primero ---; las capas
inclinadas y la \omterm{def:g7:earthquakes:fault}{falla} se desgastan hasta quedar planas; la arcilla se
deposita encima. (e) Dos: la superficie ondulada entre el esquisto y el
conglomerado, y la superficie entre las capas inclinadas y la arcilla
horizontal.
\end{solution}

\begin{solution}{exo:g7:geological-time:12}
(a) La mitad de 365 días son 182.5 días. De enero a junio hay
$31+28+31+30+31+30 = 181$ días, así que la mitad de la historia se cumple
1.5 días después de empezar julio: el \textbf{2~de julio a mediodía}.
(b) $4540 \div 2 = 2270$ millones de años. (c)
$80 \div \num{4540000000} \times \num{31536000} \approx 0.56$: unos
\textbf{0.6~segundos}.
\end{solution}

\begin{solution}{pb:g7:geological-time:1}
\textbf{1.} \num{4540} millones de años.

\textbf{2.} $4540 \div 365 \approx 12.4$ millones de años.

\textbf{3.} $\num{4540000000} \div (365 \times 24) = \num{4540000000}
\div 8760 \approx \num{518000}$ años (con el valor redondeado de \num{12.4}
millones, $\num{12400000} \div 24 \approx \num{517000}$).

\textbf{4.} $\num{4540000000} \div \num{525600} \approx 8638$: unos
\num{8600} años.

\textbf{5.} $4540 - 4030 = 510$ millones de años después de que se formara
la Tierra; $510 \div 12.4 \approx 41$ días. Enero tiene 31 días, así que
41 días han pasado al comenzar el \textbf{11~de febrero}.

\textbf{6.} $4540 - 3500 = 1040$ millones de años; $1040 \div 12.4 \approx
83.9$ días (83.6 con el valor sin redondear, $12.44$). Enero y
febrero suman 59 días: estamos entre 24.6 y 24.9 días después de empezar marzo,
el \textbf{25~de marzo}.

\textbf{7.} $66 \div 12.4 \approx 5.3$ días antes del final del año, es
decir, $365 - 5.3 = 359.7$ días después del 1~de enero. Los once primeros
meses suman 334 días: 25.7 días después de empezar diciembre, el
\textbf{26~de diciembre}.

\textbf{8.} $538.8 \div 4540 \approx 0.119$: queda el \textbf{11.9~\%} del
año, es decir, $0.119 \times 365 \approx 43$ días. $365 - 43 =
322$ días después del 1~de enero: empieza en \textbf{noviembre} (el
18~de noviembre).

\textbf{9.} $\num{10000} \div 8600 \approx 1.2$ minutos.

\textbf{10.} $5000 \div 8600 \approx 0.58$ minutos, y $0.58 \times 60
\approx 35$: unos \textbf{35 segundos} antes de medianoche.

\textbf{11.} $\num{300000} \div 8600 \approx 34.9$: unos \textbf{35
minutos} antes de medianoche.

\textbf{12.} 35 minutos antes de la medianoche del 31~de diciembre: nuestra
especie aparece hacia las \textbf{23:25}, en la última hora del año.
\end{solution}
